\SourceLeaf{8}
\pdfbookmark[0]{Disquisitio de elementis ellipticis Palladis}{palladis}
\begin{center}
{\large DISQUISITIO}\par\medskip
{\Large DE ELEMENTIS ELLIPTICIS PALLADIS}\par\medskip
EX OPPOSITIONIBUS ANNORUM 1803, 1804, 1805, 1807, 1808, 1809.
\end{center}
\Article{1}
Planetae primarii Palladis mense Martio anni 1802 detecti septem quidem hactenus fuerunt apparitiones periodicae: attamen quum planeta eo quo primum detectus est die oppositionem cum sole iam transgressus esset, sex tantummodo oppositiones hucusque numerantur; quae ad maximam partem in pluribus speculis astronomicis satis bene observatae sunt, adeoque praecisione sufficiente determinari potuerunt. Unicam oppositionem anni 1808 excipere oportet, quo anno planeta in aphelio orbitae versatus propter luminis debilitatem magnas observantibus difficultates obiecit, et ut verum fateamur, ab astronomis aliquantulum neglectus est. Ascensionum quidem rectarum copiam satis magnam suppeditavit clar. \textsc{de Lindenau}, solita cura ad tubum meridianum praestantissimum speculae Seebergensis circa tempus oppositionis observatarum: sed declinationes paucas tantum longeque minus certas aliunde obtinere potui. Quamobrem haud licuit, hancce oppositionem eadem qua ceteras praecisione stabilire, imprimisque latitudo nimis incerta mansit.
\Article{2}
Quamprimum observationes cuiusvis apparitionis finitae essent, omnes quotannis colligere, diligenter discutere, et cum observationibus annorum antecedentium combinare solitus fui, ut elementorum ellipticorum determinatio tanto ac\SourceLeaf{9}curatior inde erueretur. E postremo huiuscemodi calculo, sub finem anni 1807 instituto, elementa prodierunt in Vol. XVII 1808 Jan. Commercii literarii astronomici clar. \textsc{de Zach} publici iuris facta, omnesque quatuor oppositiones eousque observatas optime conciliantia. Anno insequente post oppositionem quintam observationum certarum penuria a novis calculis me deterruit, ad quos tunc demum reverti, ubi observationes anno 1809 circa oppositionem sextam factae computum utriusque oppositionis simul suscipere permiserunt, quem in Novis nostris literariis anni 1810 Febr. 24. indicavi.

Ecce iam conspectum omnium sex oppositionum hactenus observatarum:
\begin{center}\small\setlength{\tabcolsep}{3pt}
\begin{tabular}{l r r|r|r|r}
\multicolumn{3}{c|}{\shortstack{Tempus oppositionis ad\\meridianum Gottingensem}}&\shortstack{Dies inde ab\\initio anni 1803}&\shortstack{Longitudo\\heliocentrica}&\shortstack{Latitudo\\geocentrica}\\\hline
1803 Iun.&30.&$0^{\mathrm h}\;27^{\mathrm m}\;32^{\mathrm s}$&181.019120&$277^0\;39'\;24''0$&$+46^0\;26'\;36''0$\\
1804 Aug.&30.&$4\;58\;27$&608.207257&$337\;0\;36.1$&$+15\;1\;49.8$\\
1805 Nov.&29.&$11\;15\;4$&1064.468796&$67\;20\;42.9$&$-54\;30\;54.9$\\
1807 Maj.&4.&$14\;37\;41$&1585.609502&$223\;37\;27.7$&$+42\;11\;25.6$\\
1808 Iul.&26.&$21\;17\;32$&2034.887176&$304\;2\;59.7$&$+37\;43\;53.7$\\
1809 Sept.&22.&$16\;10\;20$&2457.673843&$359\;40\;4.4$&$-7\;22\;10.1$
\end{tabular}\end{center}
\Article{3}
Si Pallas exacte moveretur in ellipsi, atque secundum leges \textsc{Kepleri}, elementa novissima quatuor oppositiones annorum 1803, 1804, 1805, 1807 optime conciliantia levissimis tantum erroribus adhuc affecta esse possent, certoque cum oppositionibus proximis quoque intra minutum primum unum duove consentire deberent. Tantum vero abfuit a consensu tam arcto, ut potius elementa ista discrepantiam ad quatuor minuta prima ascendentem iam in oppositione quinta monstraverint, a sexta vero duodecim minutis primis aberraverint. Scilicet tantas perturbationes patitur planeta noster a reliquis praesertimque a Iove, ut consensus exactus atque stabilis inter phaenomena motumque pure ellipticum obtineri nequeat. Hinc simul patet, elementa alia atque alia proditura esse, prout oppositionibus quaternis aliis aliisve superstruantur, quod confirmatum est per calculum duorum aliorum systematum elementorum, quorum alterum ex oppositio\SourceLeaf{10}nibus annorum 1804, 1805, 1807, 1808, alterum ex oppositionibus annorum 1805, 1807, 1808, 1809 nuper deduxi. Quo melius haec diversa elementorum systemata inter se conferri possint, singula hic profero.

\Needspace{14\baselineskip}
\begin{center}I.\quad\textit{Elementa elliptica Palladis ex oppositionibus annorum}\par
1803, 1804, 1805, 1807.\end{center}
\begingroup\small\setlength{\tabcolsep}{2pt}
\noindent\begin{tabular}{@{}p{.51\linewidth}lr@{}}
Epocha longitudinis mediae pro anno 1803 ad meridianum Gottingensem&. . . . . . . . . . . . . .&$221^0\;39'\;30''4$\\
Motus diurnus medius tropicus&. . . . . . . . . . . .&$770''2143$\\
Longitudo perihelii 1803&. . . . . . . . . . . . .&$121^0\;3'\;11''4$\\
Longitudo nodi ascendentis 1803&. . . . . . . . . . .&$172\;28\;56.9$\\
Inclinatio orbitae&. . . . . . . . . . . . . . . .&$34\;37\;41.0$\\
Excentricitas $(=\sin.\,14^0\;10'\;58''81)$&. . . . . . . . . .&$0.2450198$\\
Logarithmus semiaxis maioris&. . . . . . . . . . . . .&$0.4423149$
\end{tabular}\par\endgroup

\Needspace{13\baselineskip}
\begin{center}II.\quad\textit{Elementa elliptica Palladis ex oppositionibus annorum}\par
1804, 1805, 1807, 1808.\end{center}
\begingroup\small\setlength{\tabcolsep}{2pt}
\noindent\begin{tabular}{@{}p{.51\linewidth}lr@{}}
Epocha longitudinis mediae 1803&. . . . . . . . . . .&$221^0\;34'\;56''7$\\
Motus medius tropicus diurnus&. . . . . . . . . . . .&$770''4467$\\
Longitudo perihelii 1803&. . . . . . . . . . . . .&$121^0\;5'\;22''1$\\
Longitudo nodi ascendentis 1803&. . . . . . . . . . .&$172\;28\;46.8$\\
Inclinatio orbitae&. . . . . . . . . . . . . . . .&$34\;37\;31.5$\\
Excentricitas $(=\sin.\,14^0\;10'\;4''08)$&. . . . . . . . . . .&$0.2447624$\\
Logarithmus semiaxis maioris&. . . . . . . . . . . . .&$0.4422276$
\end{tabular}\par\endgroup

\Needspace{13\baselineskip}
\begin{center}III.\quad\textit{Elementa elliptica Palladis ex oppositionibus annorum}\par
1805, 1807, 1808, 1809.\end{center}
\begingroup\small\setlength{\tabcolsep}{2pt}
\noindent\begin{tabular}{@{}p{.51\linewidth}lr@{}}
Epocha longitudinis mediae 1803&. . . . . . . . . . .&$221^0\;23'\;24''6$\\
Motus medius tropicus diurnus&. . . . . . . . . . . .&$770''9265$\\
Longitudo perihelii 1803&. . . . . . . . . . . . .&$120^0\;58'\;4''8$\\
Longitudo nodi ascendentis 1803&. . . . . . . . . . .&$172\;27\;52.4$\\
Inclinatio orbitae&. . . . . . . . . . . . . . . .&$34\;36\;49.4$\\
Excentricitas $(=\sin.\,14^0\;9'\;36''63)$&. . . . . . . . . . .&$0.2446335$\\
Logarithmus semiaxis maioris&. . . . . . . . . . . . .&$0.4420473$
\end{tabular}\par\endgroup
\SourceLeaf{11}
\Article{4}
Quonam modo orbita elliptica planetae e quatuor observationibus (quarum duae tantum completae sunt) determinari possit, in Sect. II. Libri II. Theoriae motus corporum coelestium monstravi. Solutio illic tradita problema maxima generalitate complectitur: attamen pro casu speciali, ubi quatuor observationes sunt oppositiones, methodum aliam in usum vocare praestat, cuius expositionem astronomis haud ingratam fore spero. Postulat quidem haecce methodus cognitionem approximatam singulorum elementorum, vel saltem horum quatuor, inclinationis orbitae, longitudinis nodi, longitudinis perihelii atque excentricitatis; sed nihil impedit, quominus hanc cognitionem per calculum anteriorem iam adesse supponamus. Rei summa pendet a solutione problematis sequentis:
\begin{quote}\itshape
‘Datis quatuor longitudinibus planetae in orbita, temporibus datis respondentibus, invenire longitudinem perihelii, excentricitatem, motum medium diurnum atque epocham longitudinis mediae’
\end{quote}
quam in art. sequ. tradam.
\Article{5}
Quum longitudo perihelii proxime iam nota esse supponatur, oriantur ex huius subtractione a longitudinibus in orbita anomaliae verae approximatae $v$, $v'$, $v''$, $v'''$, ubi longitudo illa vel tamquam constans spectari potest, vel si ipsius variatio annua nota fuerit, ad singula tempora rite transferenda est. Sit $e=\sin\varphi$ excentricitas approximata, qua adhibita computentur per methodos notas anomaliae mediae $M$, $M'$, $M''$, $M'''$ veris $v$, $v'$, $v''$, $v'''$ respondentes. Designando iam tempora per $t$, $t'$, $t''$, $t'''$, horum differentiae $t'-t$, $t''-t'$, $t'''-t''$ differentiis $M'-M$, $M''-M'$, $M'''-M''$ proportionales esse deberent, si a valoribus veris longitudinis perihelii atque anguli $\varphi$ profecti essemus: sin secus evenit, correctiones his quantitatibus applicandae sequenti modo eruentur. Statuendo
\[
\begin{aligned}
m&=-\frac{\cos\varphi^3}{(1+e\cos v)^2},&n&=\frac{m(2+e\cos v)\sin v}{\cos\varphi'}\\[6pt]
m'&=-\frac{\cos\varphi^3}{(1+e\cos v')^2},&n'&=\frac{m'(2+e\cos v')\sin v'}{\cos\varphi}\\[6pt]
m''&=-\frac{\cos\varphi^3}{(1+e\cos v'')^2},&n''&=\frac{m''(2+e\cos v'')\sin v''}{\cos\varphi}\\[6pt]
m'''&=-\frac{\cos\varphi^3}{(1+e\cos v''')^2},&n'''&=\frac{m'''(2+e\cos v''')\sin v'''}{\cos\varphi}
\end{aligned}
\]
